% Part: history
% Chapter: biographies
% Section: georg-cantor

\documentclass[../../../include/open-logic-section]{subfiles}

\begin{document}

\olfileid{his}{bio}{can}

\olsection{گئورگ کانتور}

\olphoto{cantor-georg}{Georg Cantor}

یکی از زندگی‌نامه‌های اولیهٔ گئورگ کانتور
(\textsc{gay}-org \textsc{kahn}-tor) مدعی بود او در کشتی‌ای که به سوی
سن پترزبورگ روسیه می‌رفت به دنیا آمد و پیدا شد و والدینش ناشناس بودند.
اما این روایت درست نیست: کانتور در سال 1845 در سن پترزبورگ به دنیا آمد
و والدینش شناخته‌شده بودند.

کانتور در سال 1867 دکترای ریاضی خود را از دانشگاه برلین گرفت. او به
سبب کارش در نظریهٔ مجموعه‌ها شناخته می‌شود و بنیان‌گذاری نظریهٔ
مجموعه‌ها به‌عنوان حوزه‌ای پژوهشی و متمایز را به او نسبت می‌دهند. او
نخستین کسی بود که اثبات کرد مجموعه‌های نامتناهی با اندازه‌های متفاوت
وجود دارند. نظریه‌های او، به‌ویژه نظریه‌اش دربارهٔ نامتناهی‌ها، در آن
زمان بحث‌های بسیاری میان ریاضی‌دانان برانگیخت و کارش مناقشه‌برانگیز بود.

باورهای دینی کانتور و کار ریاضی او جدایی‌ناپذیر بودند؛ او حتی ادعا کرد
نظریهٔ اعداد ترامتناهی را خدا مستقیماً به او الهام کرده است. کانتور در
سال‌های پایانی زندگی به بیماری روانی مبتلا شد. از سال 1894 بستری‌شدن
او آغاز شد و هرچه به پایان عمر نزدیک‌تر می‌شد دفعات آن بیشتر می‌شد.
انتقادهای سنگین از کارش، از جمله قطع رابطه با لئوپولد کرونکرِ
ریاضی‌دان، به افسردگی و بی‌علاقگی او به ریاضیات انجامید. کانتور در
دوره‌های افسردگی به فلسفه و ادبیات روی می‌آورد و حتی نظریه‌ای منتشر
کرد که فرانسیس بیکن را نویسندهٔ نمایش‌نامه‌های شکسپیر می‌دانست.

کانتور در 6 ژانویهٔ 1918 در آسایشگاهی در هاله درگذشت.

\begin{reading}
برای زندگی‌نامه‌های کامل کانتور، \citet{Dauben1990} و
\citet{Grattan-Guinness1971} را ببینید. دیدگاه‌های رادیکال کانتور در
برنامهٔ رادیویی BBC Radio 4 با عنوان
\emph{A Brief History of Mathematics} \citep{Sautoy2014} نیز توصیف
شده‌اند. اگر دوست دارید نظریه‌های کانتور را در قالب رپ بشنوید،
\citet{Rose2012} را ببینید.
\end{reading}

\end{document}
